\cvsection{開源軟體}

% \url{http://unetbootin.github.io/}
\textbf{UNetbootin}（Live USB 製作工具） \hfill \url{https://zh.wikipedia.org/wiki/UNetbootin}\\ %\hfill \textsl{2007 年 1 月 -- 至今}\\
\textcolor{sectcol}{下載量達 4,000 萬次。}UNetbootin 可為 50 多種 Linux 發行版製作可開機 USB 隨身碟。\\ %此專案已被收錄至 Debian、Ubuntu、Fedora、openSUSE、Gentoo 等主要發行版的官方套件庫。\\
% \hspace{3mm}

\vspace{-1mm}

\textbf{Wubi}（Windows 版 Ubuntu 安裝程式） \hfill \url{https://zh.wikipedia.org/wiki/Wubi}\\ %\hfill \textsl{2006 年 11 月 -- 2007 年 8 月}\\
\textcolor{sectcol}{現已成為 Ubuntu 的一部分。}開發了 Wubi 的最初幾個版本，讓使用者可從 Windows 直接安裝 Ubuntu。\\ % 此專案現已成為 Ubuntu 的一部分。\\

\vspace{-1mm}

\textbf{HabitLab}（實境行為改變研究平台） \hfill \url{https://habitlab.stanford.edu}\\ %\hfill \textsl{2006 年 11 月 -- 2007 年 8 月}\\
\textcolor{sectcol}{12,000 多名每日活躍使用者。}HabitLab 是我在博士班期間所開發，目前仍持續被史丹佛大學醫學院用於研究。\\ % 此專案現已成為 Ubuntu 的一部分。\\

% 曾任 \textbf{CHI}（2015、2018、2019、2021）、\textbf{UIST}（2017、2018）、\textbf{CSCW}（2021）、\textbf{IMWUT}（2019）審稿人。\\
